%======================================================================
\subsection{BW-ROU-06 --- Built, 30 Oct: the harness as a vault item, stock material, the purchasing BOM}\label{sec:rou06-built}
%======================================================================

Step~2 of the routing plan (\S\ref{sec:rou-next-plan}) with your answers V1--V3, built as \file{bruceworx\_harness.py},
\file{bruceworx\_stock.py}, \file{bruceworx/stock.py} (FreeCAD-free maths) and changes to \file{bruceworx\_routing.py},
\file{bruceworx\_bom.py}, \file{bruceworx/bom.py} and the Routing tab.

\paragraph{What you can do now.}
\begin{itemize}[nosep]
  \item \textbf{Harness} (Routing tab, new \emph{Harness} group) groups the selected wires, tubes, hoses, bundles (with
        their route paths) and connector links in an \code{App::Part} with \code{HarnessType} (\emph{Cable assembly}, or
        \emph{Tube assembly} when only tubes and hoses are selected) inside the assembly. The vault BOM walker reads
        through it, so the parent's BOM already lists the stock material by length and the terminals.
  \item \textbf{Make vault item} saves the container as its own file under a new P/N (vault assembly, category
        \emph{Harnesses}, the document type named like the \code{HarnessType} when the vault has one), records its BOM,
        and replaces the container in the parent by an \code{App::Link} to it. The routed objects are copied; a path
        point attached to a part of the parent becomes a cross-document link (\code{XAttachments},
        \code{App::PropertyXLinkSubList}). With both files open, moving a part re-routes the harness; opened alone
        (parent file gone) the harness keeps its last solved points and still totals itself.
  \item \textbf{Stock material.} Every wire, tube and hose panel has a \emph{Stock material} row: pick the vault item
        it is cut from or \emph{New\ldots}, which makes one from the item's properties ("Wire, 22 AWG, red, PVC",
        "Tube, 8 x 1, Nylon", "Hose, 6 x 10") asking for Manufacturer, MPN, scrap (default 10\,\%) and purchase unit
        (optional, e.g.\ 30\,m). It is a COTS vault item, category \emph{Stock material}, \code{QuantityUnit} m; an
        item with the same key is found again, not numbered twice. The routed item stores the P/N, manufacturer and MPN.
  \item \textbf{Purchasing BOM.} The Bill of Materials dialog has a \emph{Purchasing} view (and CSV): the flattened
        metered lines, summed over all harnesses and quantities, with
        $L_{\text{buy}}=\lceil (1+s)\sum n_kL_k/L_{\text{unit}}\rceil L_{\text{unit}}$ (no rounding without a purchase
        unit). Every line prints Manufacturer and MPN.
\end{itemize}

\paragraph{Decided while building.}
\begin{itemize}[nosep]
  \item Scrap and purchase unit live in the stock item's file (read from its \code{Document.xml}, no FreeCAD needed),
        not in the journal, because the vault BOM line has no field for them and \code{pdm.py} was not mine to extend.
  \item The harness is made in two steps: the file is staged with snapshot points, then checked out into the
        workspace and given its cross-document links there, because FreeCAD stores a link relative to the file it is
        saved in. A harness therefore has two versions (01 snapshot, 02 linked) until the first real check-in.
  \item New paths use \code{App::PropertyLinkSubListGlobal} for \code{Attachments}, so a path in a Harness container
        may follow geometry outside it (a local-scope link gave an ``out of the allowed scope'' warning).
\end{itemize}

\paragraph{Tests.} \file{addons/BruceWorx/tests/test\_bom.py} (scrap and rounding maths, sums over harnesses and
quantities, CSV with Manufacturer and MPN, \code{Document.xml} reader); headless \file{TestBruceWorxHarness.py} on a
scratch vault (container; Make vault item; a moved part re-routes the harness; the harness opened alone; two harnesses
of ``22 AWG red PVC'' with 10\,\% scrap and a 30\,m spool); GUI \file{TestBruceWorxRouting.py} (Harness command, Stock
row). Routing, Rope, Bom, Workspaces and the slot lint pass. \code{TestBruceWorxFlexSub} fails on
\code{test\_flexible\_copies\_and\_follows} with and without these changes (a rope has a fixed cut length, so its
\code{Length} no longer changes when a part moves); I did not touch it.

\paragraph{Open questions.}
\begin{itemize}[nosep]
  \item Parent and harness link each other, a cycle between documents: FreeCAD logs \emph{``getDependentDocuments: the
        graph must be a DAG''} on closing, though recompute and save worked. Is that acceptable, or should the parent
        hold the harness only through the vault BOM?
  \item The harness link must stay at the parent's origin (its paths are in the parent's frame). Should it be locked?
  \item A path attached to a bare App::Link's origin (sub empty) is not resolved (Links have no global placement); pin
        frames and part sub-elements are. Worth fixing in \code{attachment\_point}.
  \item Connector links in a harness keep their library file links only when that file is open at Make vault item.
\end{itemize}
